\documentclass{article}
\usepackage[utf8]{inputenc}
\usepackage[spanish]{babel}
\usepackage{amsmath, amssymb, amsthm}
\usepackage{xcolor}
\usepackage{geometry}
\usepackage{enumerate}
\usepackage{mdframed}
\geometry{a4paper, margin=2.5cm}

\newtheorem*{teo}{Teorema}
\newtheorem*{defn}{Definición}
\newtheorem*{cor}{Corolario}
\newtheorem*{prop}{Proposición}
\newtheorem*{lema}{Lema}
\newenvironment{marco}{\begin{mdframed}[linecolor=black!40,backgroundcolor=gray!5]\small\textbf{Marco teórico.}\par\smallskip}{\end{mdframed}\medskip}

\title{\textbf{Guía 1 --- EII-801 Programación Matemática\\[4pt]
\large Preguntas y Soluciones con Desarrollo Completo\\[2pt]
\normalsize Referencias: Bertsimas \& Tsitsiklis (BT); Chvátal (Ch)}}
\author{Pitehr Hurtado}
\date{Primer semestre 2026}

\begin{document}

\maketitle
\tableofcontents
\newpage

% ============================================================
\section*{Problema 1 --- Diccionario simplex con parámetros}
\addcontentsline{toc}{section}{Problema 1}

\textbf{Enunciado.} Se tiene el siguiente diccionario de un PL en forma estándar (maximización):
\begin{align*}
X_3 &= c - 4X_1 + aX_2 + bX_5 \\
X_4 &= d + X_1 + 5X_2 + X_5 \\
X_6 &= 3 + eX_1 + 3X_2 + 4X_5 \\
Z   &= 80 + fX_1 + gX_2 + hX_5
\end{align*}
Proponga valores para los parámetros $a,b,c,d,e,f,g,h$ en cada uno de los casos:
\begin{enumerate}[a)]
  \item diccionario óptimo con soluciones múltiples,
  \item solución básica factible (SBF) degenerada,
  \item SBF no degenerada y no óptima; realizar una iteración del simplex,
  \item escribir la formulación original del problema del caso (c).
\end{enumerate}

\begin{marco}
\textbf{[BT §3.1]} En un diccionario, las variables no básicas valen cero; los valores de la SBF son los términos independientes de cada fila. Factibilidad: todos los términos independientes $\geq 0$. \\[2pt]
\textbf{[BT §3.2]} Criterio de optimalidad (maximización): todos los costos reducidos de las variables no básicas $\leq 0$. \\[2pt]
\textbf{[BT §3.1]} Degeneración: alguna variable básica toma valor $0$ en la SBF. \\[2pt]
\textbf{[BT §3.3]} Soluciones múltiples óptimas: existe algún costo reducido igual a $0$ en el óptimo; la dirección correspondiente es factible ($\theta^* > 0$).
\end{marco}

\subsection*{Solución}

Las variables no básicas son $\{X_1, X_2, X_5\} = 0$. La SBF actual es:
$X_3 = c,\; X_4 = d,\; X_6 = 3,\; Z = 80$.

\textbf{Condiciones generales:}
\begin{itemize}
  \item Factibilidad: $c \geq 0$, $d \geq 0$, $3 \geq 0$ (automático).
  \item Optimalidad (max): $f \leq 0$, $g \leq 0$, $h \leq 0$.
  \item Degeneración: $c = 0$ o $d = 0$.
\end{itemize}

\begin{enumerate}[a)]

\item \textbf{Óptimo con soluciones múltiples.}

Tomemos $c=1,\; d=2,\; a=2,\; b=2,\; e=1,\; f=0,\; g=-1,\; h=-1$.

El diccionario queda:
\begin{align*}
X_3 &= 1 - 4X_1 + 2X_2 + 2X_5 \\
X_4 &= 2 + X_1 + 5X_2 + X_5 \\
X_6 &= 3 + X_1 + 3X_2 + 4X_5 \\
Z   &= 80 + 0 \cdot X_1 - X_2 - X_5
\end{align*}

Verificación \textbf{[BT §3.2]}:
\begin{itemize}
  \item Factibilidad: $X_3=1>0$, $X_4=2>0$, $X_6=3>0$ \checkmark
  \item Optimalidad: $f=0\leq 0$, $g=-1\leq 0$, $h=-1\leq 0$ \checkmark
  \item Soluciones múltiples: $f=0$, por lo que $X_1$ puede entrar sin cambiar $Z$. Razón mínima: solo la fila de $X_3$ tiene coeficiente $-4 < 0$ para $X_1$, así que $\theta^* = 1/4 > 0$ \checkmark. La dirección es factible, confirmando múltiples óptimos \textbf{[BT §3.3]}.
\end{itemize}

\item \textbf{SBF degenerada.}

Tomemos $c=0,\; d=2,\; a=1,\; b=1,\; e=1,\; f=1,\; g=-1,\; h=-1$.

\begin{itemize}
  \item $X_3 = 0$ (degeneración \textbf{[BT §3.1]}), $X_4=2>0$, $X_6=3>0$.
  \item $f=1>0$: no es óptimo (hay margen de mejora). Sirve como ejemplo de SBF degenerada no óptima.
\end{itemize}

\item \textbf{SBF no degenerada y no óptima; una iteración.}

Tomemos $c=2,\; d=3,\; a=1,\; b=1,\; e=1,\; f=2,\; g=-1,\; h=-1$.

Diccionario inicial:
\begin{align*}
X_3 &= 2 - 4X_1 + X_2 + X_5 \\
X_4 &= 3 + X_1 + 5X_2 + X_5 \\
X_6 &= 3 + X_1 + 3X_2 + 4X_5 \\
Z   &= 80 + 2X_1 - X_2 - X_5
\end{align*}

No degenerada: $X_3=2>0$, $X_4=3>0$, $X_6=3>0$. No óptima: $f=2>0$ \textbf{[BT §3.2]}.

\textit{Iteración del simplex:} Entra $X_1$ (el único costo reducido positivo). Calculamos las razones para las filas con coeficiente negativo de $X_1$:
\[
  \theta^* = \min\left\{\frac{2}{4}\right\} = \frac{1}{2}, \quad \text{sale } X_3.
\]
(Las filas $X_4$ y $X_6$ tienen coeficiente de $X_1$ positivo, por lo que no restringen.)

\textbf{Pivote:} De la fila de $X_3$:
\[
  X_1 = \frac{1}{2} + \frac{1}{4}X_2 + \frac{1}{4}X_5 - \frac{1}{4}X_3.
\]

Sustituyendo en las demás filas:
\begin{align*}
X_4 &= 3 + \left(\frac{1}{2}+\frac{1}{4}X_2+\frac{1}{4}X_5-\frac{1}{4}X_3\right) + 5X_2 + X_5 \\
    &= \frac{7}{2} + \frac{21}{4}X_2 + \frac{5}{4}X_5 - \frac{1}{4}X_3 \\[4pt]
X_6 &= 3 + \left(\frac{1}{2}+\frac{1}{4}X_2+\frac{1}{4}X_5-\frac{1}{4}X_3\right) + 3X_2 + 4X_5 \\
    &= \frac{7}{2} + \frac{13}{4}X_2 + \frac{17}{4}X_5 - \frac{1}{4}X_3 \\[4pt]
Z   &= 80 + 2\left(\frac{1}{2}+\frac{1}{4}X_2+\frac{1}{4}X_5-\frac{1}{4}X_3\right) - X_2 - X_5 \\
    &= 81 - \frac{1}{2}X_2 - \frac{1}{2}X_5 - \frac{1}{2}X_3
\end{align*}

Nuevo diccionario:
\begin{align*}
X_1 &= \frac{1}{2} + \frac{1}{4}X_2 + \frac{1}{4}X_5 - \frac{1}{4}X_3 \\
X_4 &= \frac{7}{2} + \frac{21}{4}X_2 + \frac{5}{4}X_5 - \frac{1}{4}X_3 \\
X_6 &= \frac{7}{2} + \frac{13}{4}X_2 + \frac{17}{4}X_5 - \frac{1}{4}X_3 \\
Z   &= 81 - \frac{1}{2}X_2 - \frac{1}{2}X_5 - \frac{1}{2}X_3
\end{align*}

Todos los costos reducidos son $-\frac{1}{2} < 0$ \textbf{[BT §3.2]}: óptimo alcanzado. $Z^* = 81$, SBF $= \left(\tfrac{1}{2}, \tfrac{7}{2}, \tfrac{7}{2}\right)$ no degenerada.

\item \textbf{Formulación original del problema (c).}

Las holguras son $X_3, X_4, X_6 \geq 0$. Leyendo el diccionario inicial:
\begin{itemize}
  \item Fila $X_3$: $X_3 = 2 - 4X_1 + X_2 + X_5$ corresponde a $4X_1 - X_2 - X_5 \leq 2$.
  \item Fila $X_4$: $X_4 = 3 + X_1 + 5X_2 + X_5$ corresponde a $-X_1 - 5X_2 - X_5 \leq 3$.
  \item Fila $X_6$: $X_6 = 3 + X_1 + 3X_2 + 4X_5$ corresponde a $-X_1 - 3X_2 - 4X_5 \leq 3$.
  \item Objetivo: $\max Z = 2X_1 - X_2 - X_5$.
\end{itemize}

\[
\begin{array}{ll}
\max & 2X_1 - X_2 - X_5 \\[2pt]
\text{s.a.} & 4X_1 - X_2 - X_5 \leq 2 \\
& -X_1 - 5X_2 - X_5 \leq 3 \\
& -X_1 - 3X_2 - 4X_5 \leq 3 \\
& X_1, X_2, X_5 \geq 0
\end{array}
\]

\end{enumerate}

% ============================================================
\section*{Problema 2 --- Diccionario óptimo de producción}
\addcontentsline{toc}{section}{Problema 2}

\textbf{Enunciado.} Se tiene el siguiente diccionario óptimo:
\begin{align*}
X_4 &= 4 + 3X_1 + X_2 + 2X_5 \\
X_3 &= 4 - 2X_1 - 2X_2 - X_5 \\
Z   &= 28 - 12X_1 - 2X_2 - 7X_5
\end{align*}
\begin{enumerate}[a)]
  \item ¿Cuántos productos y recursos tiene el problema original?
  \item Escribir la formulación original.
  \item Escribir el dual y su diccionario óptimo.
  \item ¿Conviene expandir el recurso 2 a un costo fijo de \$3000 más \$3/unidad adicional?
\end{enumerate}

\begin{marco}
\textbf{[BT §3.2]} Variables básicas: $\{X_3, X_4\}$; no básicas: $\{X_1, X_2, X_5\} = 0$. SBF óptima: $X_3=4$, $X_4=4$, $Z=28$. \\[2pt]
\textbf{[BT §4.4]} Los precios sombra (duales) se leen del diccionario óptimo: $y_i^* = $ coeficiente de la holgura $s_i$ en la fila del objetivo del diccionario óptimo (con signo adecuado). \\[2pt]
\textbf{[BT §4.3 Teo.\ 4.5]} Holgura complementaria: si $x_j^* > 0$ entonces la $j$-ésima restricción dual es activa; si $y_i^* > 0$ entonces la $i$-ésima restricción primal es activa. \\[2pt]
\textbf{[BT §5.3]} El precio sombra $y_i^*$ representa el cambio en $Z^*$ por unidad de incremento en $b_i$, válido en el rango donde la misma base permanece factible.
\end{marco}

\subsection*{Solución}

\begin{enumerate}[a)]

\item \textbf{Productos y recursos.}

Variables básicas: $\{X_3, X_4\}$ $\Rightarrow$ $m = 2$ restricciones (recursos). Variables no básicas de decisión: $\{X_1, X_2\}$; la variable $X_5$ es la holgura del recurso 2 (que sale de la base al hacerse activa). Así el problema tiene \textbf{2 productos} ($X_1, X_2$) y \textbf{2 recursos} (holguras $X_4, X_5$ en la base inicial). $X_3$ entró a la base durante la resolución.

En términos de producción: 3 variables de decisión ($X_1, X_2, X_3$) y 2 recursos, donde $X_4$ y $X_5$ son las holguras. En el diccionario óptimo $X_3$ es básica (se fabrica), mientras $X_1$ y $X_2$ son no básicas (no se fabrican).

\item \textbf{Formulación original.}

Del diccionario óptimo se reconstruye el problema invirtiendo las operaciones de pivote. La fila de $X_3$:
\[
  X_3 = 4 - 2X_1 - 2X_2 - X_5 \;\Rightarrow\; 2X_1 + 2X_2 + X_3 + X_5 = 4
\]
(restricción del recurso 2: $2X_1 + 2X_2 + X_3 \leq 4$, con holgura $X_5$).

La fila de $X_4$:
\[
  X_4 = 4 + 3X_1 + X_2 + 2X_5
\]
Dado que $X_5$ es la holgura del recurso 2, necesitamos expresar $X_4$ en términos originales. Sustituyendo $X_5 = 4 - 2X_1 - 2X_2 - X_3$:
\[
  X_4 = 4 + 3X_1 + X_2 + 2(4 - 2X_1 - 2X_2 - X_3) = 12 - X_1 - 3X_2 - 2X_3
\]
Luego $X_4 = 12 - X_1 - 3X_2 - 2X_3$, es decir, la restricción del recurso 1 es $X_1 + 3X_2 + 2X_3 \leq 12$.

Para el objetivo: en el diccionario óptimo $Z = 28 - 12X_1 - 2X_2 - 7X_5$. Con $c_3$ siendo el coeficiente de $X_3$ (básica), el precio sombra de $X_3$ implica: a partir de la fila de $X_3$, el coeficiente del objetivo asociado a $X_3$ es $c_3 = 7$ (valor del recurso que aporta). Reconstruyendo: $Z = c_1 X_1 + c_2 X_2 + c_3 X_3$ con $c_3 = 7$, $c_1 = 2$, $c_2 = 12$ (se verifica que $Z^* = 0 + 0 + 7(4) = 28$ \checkmark).

\medskip
\textbf{Problema original:}
\[
\begin{array}{ll}
\max & 2X_1 + 12X_2 + 7X_3 \\
\text{s.a.} & X_1 + 3X_2 + 2X_3 \leq 12 \quad \text{(recurso 1)} \\
& 2X_1 + 2X_2 + X_3 \leq 4 \quad \text{(recurso 2)} \\
& X_1, X_2, X_3 \geq 0
\end{array}
\]

\textit{Verificación:} $(X_1^*, X_2^*, X_3^*) = (0,0,4)$: recurso 1: $2(4)=8\leq 12$ \checkmark; recurso 2: $4 \leq 4$ \checkmark; $Z=28$ \checkmark.

\item \textbf{Dual y su diccionario óptimo.}

\[
\begin{array}{ll}
\min & W = 12y_1 + 4y_2 \\
\text{s.a.} & y_1 + 2y_2 \geq 2 \\
& 3y_1 + 2y_2 \geq 12 \\
& 2y_1 + y_2 \geq 7 \\
& y_1, y_2 \geq 0
\end{array}
\]

Por holgura complementaria \textbf{[BT §4.3 Teo.\ 4.5]}: $X_4^* = 4 > 0 \Rightarrow y_1^* = 0$; $X_3^* = 4 > 0 \Rightarrow$ restricción dual 3 activa: $2(0) + y_2^* = 7 \Rightarrow y_2^* = 7$.

Verificación: $y_1^* + 2y_2^* = 14 \geq 2$ \checkmark; $3(0)+2(7)=14\geq 12$ \checkmark; $W^* = 0+4(7) = 28 = Z^*$ \checkmark (dualidad fuerte \textbf{[BT §4.2 Teo.\ 4.2]}).

\medskip
Diccionario óptimo del dual (no básicas $\{y_1, e_3\} = 0$, donde $e_j$ es la holgura de la restricción dual $j$):
\begin{align*}
y_2    &= 7 - 2y_1 + e_3 \\
e_1    &= 12 - 3y_1 + 2e_3 \\
e_2    &= 2 - y_1 + 2e_3 \\
W      &= 28 + 4y_1 + 4e_3
\end{align*}

Dado que es minimización, los costos reducidos de las no básicas deben ser $\geq 0$: ambos valen $+4 > 0$ \textbf{[BT §3.2]} (convenio para min). Óptimo confirmado.

\item \textbf{Conveniencia de expandir el recurso 2.}

Precio sombra del recurso 2: $y_2^* = 7$ \$/unidad \textbf{[BT §5.3]}. La ganancia bruta por expandir $\Delta$ unidades es $7\Delta$. El costo es $3000 + 3\Delta$.

Ganancia neta: $7\Delta - 3000 - 3\Delta = 4\Delta - 3000$.

Requiere $\Delta > 750$ para ser rentable. Ahora verificamos el rango de validez del precio sombra. La columna de $X_5$ (holgura del recurso 2) en $B^{-1}$ actúa sobre el vector básico:
\[
  X_3' = 4 - \Delta \geq 0 \;\Rightarrow\; \Delta \leq 4.
\]
(Leer de la fila de $X_3$: coeficiente de $X_5$ es $-1$.)

Dentro del rango $[0,4]$: ganancia neta máxima $= 4(4) - 3000 = 16 - 3000 = -2984 < 0$.

\textbf{Conclusión: No conviene expandir el recurso 2.} Para cualquier $\Delta$ dentro del rango de validez del precio sombra, la expansión es deficitaria.

\end{enumerate}

% ============================================================
\section*{Problema 3 --- Problema agrícola (sensibilidad)}
\addcontentsline{toc}{section}{Problema 3}

\textbf{Enunciado.} Problema de producción agrícola:
\[
\begin{array}{ll}
\max & 150x_1 + 200x_2 - 10x_3 \\
\text{s.a.} & x_1 + x_2 \leq 45 \quad (1) \\
& 6x_1 + 10x_2 - x_3 \leq 0 \quad (2) \\
& x_3 \leq 350 \quad (3) \\
& 5x_1 \leq 140 \quad (4) \\
& 4x_2 \leq 120 \quad (5) \\
& x_1, x_2, x_3 \geq 0
\end{array}
\]
Un solver entrega: $Z^* = 4250$, $(x_1^*, x_2^*, x_3^*) = (25, 20, 350)$. Variables duales: $y_1^* = 75$, $y_2^* = 12.5$, $y_3^* = 2.5$, $y_4^* = y_5^* = 0$.
\begin{enumerate}[a)]
  \item Escribir el diccionario óptimo.
  \item Precio máximo que se justifica pagar por 1 ha adicional adyacente.
  \item El precio del trigo baja, haciendo que el ingreso de $x_1$ baje a 130 M\$/ha.
  \item Se inundan 5 ha: nueva solución.
  \item Se inundan 8 ha: nueva solución.
\end{enumerate}

\begin{marco}
\textbf{[BT §4.4 / §5.3]} Precio sombra $y_i^* = \partial Z^*/\partial b_i$: valor marginal del recurso $i$. Válido para $\Delta b_i$ en el rango donde la misma base óptima permanece factible. \\[2pt]
\textbf{[BT §4.3 Teo.\ 4.5]} $y_i^* > 0 \Leftrightarrow$ restricción $i$ activa en el óptimo. \\[2pt]
\textbf{[BT §5.1]} Cambio en $c_j$ con $x_j$ básica: la SBF no cambia; se recalculan los precios duales y los costos reducidos. \\[2pt]
\textbf{[BT §5.3]} Al cambiar $b$ en $\Delta e_i$: $x_B' = x_B + \Delta \cdot B^{-1}e_i$. La base sigue siendo factible ssi $x_B' \geq 0$.
\end{marco}

\subsection*{Solución}

\begin{enumerate}[a)]

\item \textbf{Diccionario óptimo.}

Restricciones activas: $(1): x_1+x_2=45$, $(2): 6x_1+10x_2=x_3$, $(3): x_3=350$. Sus holguras $s_1=s_2=s_3=0$ son no básicas. Restricciones inactivas: $s_4 = 140-5(25)=15>0$, $s_5=120-4(20)=40>0$; $s_4$ y $s_5$ son básicas. Base óptima: $\mathcal{B} = \{x_1, x_2, x_3, s_4, s_5\}$.

El diccionario óptimo expresa las variables básicas en función de $\{s_1, s_2, s_3\}$. Para construirlo, escribimos el sistema $B x_{\mathcal{B}} = b - N x_N$ y resolvemos $B^{-1}$. El resultado es:

\begin{align*}
x_1  &= 25 - \tfrac{5}{2}s_1 + \tfrac{1}{4}s_2 + \tfrac{1}{4}s_3 \\
x_2  &= 20 + \tfrac{3}{2}s_1 - \tfrac{1}{4}s_2 - \tfrac{1}{4}s_3 \\
x_3  &= 350 - s_3 \\
s_4  &= 15 + \tfrac{25}{2}s_1 - \tfrac{5}{4}s_2 - \tfrac{5}{4}s_3 \\
s_5  &= 40 - 6s_1 + s_2 + s_3 \\
Z    &= 4250 - 75s_1 - 12.5\,s_2 - 2.5\,s_3
\end{align*}

Los costos reducidos de todas las no básicas son negativos (o cero), confirmando optimalidad \textbf{[BT §3.2]}.

\item \textbf{Precio máximo por 1 ha adyacente.}

El precio sombra de la restricción (1) es $y_1^* = 75$ M\$/ha \textbf{[BT §4.4]}. Adquirir 1 ha adicional incrementa $Z^*$ en 75 M\$. El propietario puede cobrar \textbf{máximo 75 M\$} por esa hectárea adicional.

\item \textbf{Baja de ingreso de $x_1$: nuevo $c_1' = 130$.}

$x_1$ es básica, por lo que la SBF $\{25, 20, 350\}$ no cambia \textbf{[BT §5.1]}. Se recalculan los precios duales $\mu' = c_{\mathcal{B}}'B^{-1}$. Solo cambia el coeficiente de $x_1$: $c_1' = 130$ (antes 150).

Usando los coeficientes de $B^{-1}$ (columnas de $s_1, s_2, s_3$ en la fila de $x_1$):
\[
  \mu_1' = 130\cdot\left(-\tfrac{5}{2}\right) + 200\cdot\left(\tfrac{3}{2}\right) + (-10)\cdot 0 = -325 + 300 = -25
\]
Pero se obtiene más directamente usando el cambio de costo: $\Delta Z = \Delta c_1 \cdot x_1^* = (-20)(25) = -500$ M\$.

Los nuevos costos reducidos de las no básicas son:
\begin{align*}
\bar{c}_{s_1}' &= 0 - 130\left(\tfrac{5}{2}\right) - 200\left(-\tfrac{3}{2}\right) + 10\cdot 0 = -325+300 = -25 \leq 0 \\
\bar{c}_{s_2}' &= 0 - 130\left(-\tfrac{1}{4}\right) - 200\left(\tfrac{1}{4}\right) + 0 = 32.5 - 50 = -17.5 \leq 0 \\
\bar{c}_{s_3}' &= 0 - 130\left(-\tfrac{1}{4}\right) - 200\left(\tfrac{1}{4}\right) + 10 = 32.5 - 50 + 10 = -7.5 \leq 0
\end{align*}

Todos los costos reducidos son $\leq 0$: la base sigue siendo óptima \textbf{[BT §5.1]}.

Nuevo $Z' = 130(25) + 200(20) - 10(350) = 3250 + 4000 - 3500 = \mathbf{3750}$ M\$. La utilidad baja en 500 M\$.

\item \textbf{5 ha inundadas: $b_1 \to 40$, $\Delta = -5$.}

Actualizar la SBF \textbf{[BT §5.3]}: $x_{\mathcal{B}}' = x_{\mathcal{B}} + \Delta \cdot B^{-1}e_1$. La columna $B^{-1}e_1$ se lee de las entradas de $s_1$ en el diccionario óptimo: $\bigl(-\tfrac{5}{2},\;\tfrac{3}{2},\;0,\;\tfrac{25}{2},\;-6\bigr)'$.

\begin{align*}
x_1' &= 25 + (-5)\left(-\tfrac{5}{2}\right) = 25 + 12.5 = 37.5 \geq 0 \;\checkmark\\
x_2' &= 20 + (-5)\left(\tfrac{3}{2}\right) = 20 - 7.5 = 12.5 \geq 0 \;\checkmark\\
x_3' &= 350 + (-5)(0) = 350 \geq 0 \;\checkmark\\
s_4' &= 15 + (-5)\left(\tfrac{25}{2}\right) = 15 - 62.5 = -47.5 < 0 \;\text{??}
\end{align*}

Espera: los signos en la columna de $s_1$ en el diccionario representan los coeficientes con los que $s_1$ aparece en cada fila. Como $s_1$ es la holgura de la restricción (1), aumentar $b_1$ equivale a aumentar $s_1$; no obstante, para analizar el cambio en $b_1$ directamente usamos $\Delta b_1 = -5$ con la columna correspondiente de $B^{-1}$.

Leyendo los coeficientes de $s_1$ en el diccionario \textbf{con signo cambiado} (pues $s_1$ aparece con coeficiente negativo cuando la restricción original es $\leq$):
La columna de $B^{-1}$ asociada a la restricción 1 es $\bigl(\tfrac{5}{2},\;{-\tfrac{3}{2}},\;0,\;{-\tfrac{25}{2}},\;6\bigr)'$.

\begin{align*}
x_1' &= 25 + (-5)\left(\tfrac{5}{2}\right) = 25 - 12.5 = 12.5 \geq 0\;\checkmark\\
x_2' &= 20 + (-5)\left(-\tfrac{3}{2}\right) = 20 + 7.5 = 27.5 \geq 0\;\checkmark\\
x_3' &= 350 \geq 0\;\checkmark\\
s_4' &= 15 + (-5)\left(-\tfrac{25}{2}\right) = 15 + 62.5 = 77.5 \geq 0\;\checkmark\\
s_5' &= 40 + (-5)(6) = 40 - 30 = 10 \geq 0\;\checkmark
\end{align*}

Base factible. Nuevo $Z' = 4250 + 75(-5) = 4250 - 375 = \mathbf{3875}$ M\$.

\item \textbf{8 ha inundadas: $b_1 \to 37$, $\Delta = -8$.}

\begin{align*}
s_5' &= 40 + (-8)(6) = 40 - 48 = -8 < 0
\end{align*}

La base se vuelve infactible \textbf{[BT §5.3]}. El rango de validez es $s_5' = 40 + 6\Delta \geq 0 \Rightarrow \Delta \geq -\tfrac{40}{6} \approx -6.67$. Con $\Delta = -8$ superamos este límite y debemos re-optimizar.

Al salirse $s_5$ de la factibilidad, la restricción (5) se hace activa: $x_2 = 30$. El nuevo problema restringido con $x_1 + x_2 \leq 37$ y $x_2 \leq 30$:

Sustituyendo $x_3 = 6x_1 + 10x_2$ (restricción 2 activa) y $x_2 = 30$:
\[
  Z = 150x_1 + 200(30) - 10(6x_1 + 10(30)) = 150x_1 + 6000 - 60x_1 - 3000 = 90x_1 + 3000.
\]

Restricciones: $x_1 \leq 37 - 30 = 7$, $5x_1 \leq 140 \Rightarrow x_1 \leq 28$.

Activo: $x_1^* = 7$. Entonces $x_2^* = 30$, $x_3^* = 6(7)+10(30) = 342$.

$Z^* = 150(7) + 200(30) - 10(342) = 1050 + 6000 - 3420 = \mathbf{3630}$ M\$. $\Delta Z = -620$ M\$.

\end{enumerate}

% ============================================================
\section*{Problema 4 --- Dual de un modelo de juego}
\addcontentsline{toc}{section}{Problema 4}

\textbf{Enunciado.} Construir el dual del siguiente problema (teoría de juegos):
\[
\begin{array}{ll}
\min & v \\
\text{s.a.} & \sum_{j=1}^{n} a_{ij} x_j \leq v, \quad i = 1,\ldots,m \\
& \sum_{j=1}^{n} x_j = 1 \\
& x_j \geq 0, \quad j=1,\ldots,n
\end{array}
\]

\begin{marco}
\textbf{[BT §4.2]} Reglas de dualización: (i) para $\min$ con restricción $\leq$: variable dual $\geq 0$; (ii) para restricción de igualdad: variable dual $\lambda$ libre; (iii) para variable primal libre: restricción dual es igualdad; (iv) para variable primal $\geq 0$ en $\min$: restricción dual $\geq 0$ (tipo $\geq$).
\end{marco}

\subsection*{Solución}

Las variables primales son $(v, x_1, \ldots, x_n)$.

\textbf{Reescritura:} las restricciones $i$ se escriben como $\sum_j a_{ij}x_j - v \leq 0$ (asociamos duals $y_i \geq 0$). La restricción de igualdad $\sum_j x_j = 1$ tiene dual $\lambda$ libre.

\textbf{Restricción dual de $v$} (variable libre en el primal):
\[
  \sum_{i=1}^{m}(-1)\cdot y_i = 1 \;\Rightarrow\; -\sum_{i=1}^m y_i = 1.
\]

\textbf{Restricción dual de $x_j$} ($x_j \geq 0$ en $\min$, restricción tipo $\geq$):
\[
  \sum_{i=1}^m a_{ij} y_i + \lambda \geq 0.
\]

\textbf{Objetivo dual:} $\max\; 0\cdot y + 1\cdot\lambda = \lambda$.

El dual queda:
\[
\begin{array}{ll}
\max & \lambda \\
\text{s.a.} & \sum_{i=1}^m a_{ij} y_i + \lambda \leq 0, \quad j=1,\ldots,n \\
& \sum_{i=1}^m y_i = -1 \\
& y_i \geq 0, \quad i=1,\ldots,m
\end{array}
\]

\textbf{Interpretación:} con el cambio $p_i = -y_i \leq 0$ y $w = -\lambda$ se obtiene la forma usual del problema de la estrategia del jugador fila en un juego de suma cero con matriz de pagos $A$. El dual maximiza el pago garantizado para el jugador columna.

% ============================================================
\section*{Problema 5 --- Planeación de producción (sensibilidad)}
\addcontentsline{toc}{section}{Problema 5}

\textbf{Enunciado.}
\[
\begin{array}{ll}
\max & x_1 + x_2 + 3x_3 + 2x_4 \\
\text{s.a.} & 2x_1 + 2x_2 + 5x_3 + 5x_4 \leq 20 \\
& 2x_1 + x_2 + 4x_3 + 2x_4 \leq 12 \\
& x_1, x_2, x_3, x_4 \geq 0
\end{array}
\]
Solución óptima: $(x_1, x_2, x_3, x_4) = (0,\, 20/3,\, 4/3,\, 0)$.
\begin{enumerate}[a)]
  \item Diccionario inicial.
  \item Diccionario óptimo sin resolver desde cero.
  \item Si $c_1$ sube a 2, ¿el nuevo costo reducido de $x_1$ sigue $\leq 0$? ¿Es aún óptima la base?
  \item Si la respuesta anterior es sí, ¿cuánto debe subir $c_1$ para que $x_1$ entre a la base?
  \item Si la respuesta anterior es no, re-optimizar.
  \item Si $a_{13}$ baja de 5 a 3, encontrar la nueva solución óptima.
\end{enumerate}

\begin{marco}
\textbf{[BT §3.1]} Diccionario inicial con holguras $s_1, s_2$. \\[2pt]
\textbf{[BT §4.4 / §5.1]} Con base $\mathcal{B} = \{x_2, x_3\}$: calcular $B^{-1}$, precios duales $\mu = c_{\mathcal{B}}B^{-1}$, costos reducidos $\bar{c}_j = c_j - \mu a_j$. \\[2pt]
\textbf{[BT §5.1]} Cambio en $c_j$ (no básica): nuevo $\bar{c}_j' = c_j' - \mu' a_j$. Si $\bar{c}_j' \leq 0$, la base sigue óptima. \\[2pt]
\textbf{[BT §5.3]} Cambio en $a_{ij}$ con $x_j$ básica: nueva base $B'$, recalcular $B'^{-1}$ y costos reducidos.
\end{marco}

\subsection*{Solución}

\begin{enumerate}[a)]

\item \textbf{Diccionario inicial.}
\begin{align*}
s_1 &= 20 - 2x_1 - 2x_2 - 5x_3 - 5x_4 \\
s_2 &= 12 - 2x_1 - x_2 - 4x_3 - 2x_4 \\
Z   &= x_1 + x_2 + 3x_3 + 2x_4
\end{align*}

\item \textbf{Diccionario óptimo (sin simplex paso a paso).}

La solución óptima $(x_2, x_3) = (20/3, 4/3)$ indica base $\mathcal{B} = \{x_2, x_3\}$ con columnas:
\[
  B = \begin{pmatrix} 2 & 5 \\ 1 & 4 \end{pmatrix}, \quad \det(B) = 8-5 = 3.
\]
\[
  B^{-1} = \frac{1}{3}\begin{pmatrix} 4 & -5 \\ -1 & 2 \end{pmatrix}.
\]

Verificación: $x_{\mathcal{B}} = B^{-1}b = \frac{1}{3}\begin{pmatrix}4(20)-5(12)\\-1(20)+2(12)\end{pmatrix} = \frac{1}{3}\begin{pmatrix}20\\4\end{pmatrix} = \begin{pmatrix}20/3\\4/3\end{pmatrix}$ \checkmark.

Precios duales: $\mu = c_{\mathcal{B}}B^{-1} = (1,3)\cdot\frac{1}{3}\begin{pmatrix}4&-5\\-1&2\end{pmatrix} = \frac{1}{3}(4-3,\;-5+6) = \left(\frac{1}{3},\frac{1}{3}\right)$.

Costos reducidos \textbf{[BT §5.1]}:
\begin{align*}
\bar{c}_{x_1} &= 1 - \left(\tfrac{1}{3},\tfrac{1}{3}\right)(2,2)' = 1 - \tfrac{4}{3} = -\tfrac{1}{3} \\
\bar{c}_{x_4} &= 2 - \left(\tfrac{1}{3},\tfrac{1}{3}\right)(5,2)' = 2 - \tfrac{7}{3} = -\tfrac{1}{3} \\
\bar{c}_{s_1} &= 0 - \left(\tfrac{1}{3},\tfrac{1}{3}\right)(1,0)' = -\tfrac{1}{3} \\
\bar{c}_{s_2} &= 0 - \left(\tfrac{1}{3},\tfrac{1}{3}\right)(0,1)' = -\tfrac{1}{3}
\end{align*}

Diccionario óptimo (invirtiendo $B$):
\begin{align*}
x_2 &= \tfrac{20}{3} + \tfrac{2}{3}x_1 - \tfrac{10}{3}x_4 - \tfrac{4}{3}s_1 + \tfrac{5}{3}s_2 \\
x_3 &= \tfrac{4}{3}  - \tfrac{2}{3}x_1 + \tfrac{1}{3}x_4 + \tfrac{1}{3}s_1 - \tfrac{2}{3}s_2 \\
Z   &= \tfrac{32}{3} - \tfrac{1}{3}x_1 - \tfrac{1}{3}x_4 - \tfrac{1}{3}s_1 - \tfrac{1}{3}s_2
\end{align*}

\item \textbf{$c_1$ sube a 2: ¿sigue siendo óptima la base?}

$x_1$ es no básica. Nuevo costo reducido \textbf{[BT §5.1]}:
\[
  \bar{c}_{x_1}' = 2 - \left(\tfrac{1}{3},\tfrac{1}{3}\right)(2,2)' = 2 - \tfrac{4}{3} = \tfrac{2}{3} > 0.
\]
\textbf{No}, la base ya no es óptima. $x_1$ tiene costo reducido positivo y debería entrar a la base.

\item \textbf{¿Cuánto debe subir $c_1$ para que $x_1$ entre?}

No aplica: según (c), ya con $c_1 = 2$ la base deja de ser óptima.

El umbral exacto es: $\bar{c}_{x_1}' = c_1' - \frac{4}{3} = 0 \Rightarrow c_1' = \frac{4}{3} \approx 1.33$. Con el valor original $c_1 = 1 < \frac{4}{3}$, la base era óptima. Cualquier $c_1' > \frac{4}{3}$ hace que $x_1$ quiera entrar.

\item \textbf{Re-optimizar con $c_1 = 2$.}

Entra $x_1$ ($\bar{c}_{x_1}' = \tfrac{2}{3} > 0$). Razón mínima: solo la fila de $x_3$ tiene coeficiente negativo para $x_1$ ($-\tfrac{2}{3}$):
\[
  \theta^* = \frac{4/3}{2/3} = 2. \quad \text{Sale } x_3.
\]

Pivote sobre la fila de $x_3$: $x_1 = 2 - \tfrac{3}{2}x_3 + \tfrac{1}{2}x_4 + \tfrac{1}{2}s_1 - s_2$.

Actualizar fila $x_2$:
\[
  x_2 = \tfrac{20}{3} + \tfrac{2}{3}\!\left(2 - \tfrac{3}{2}x_3 + \tfrac{1}{2}x_4 + \tfrac{1}{2}s_1 - s_2\right) - \tfrac{10}{3}x_4 - \tfrac{4}{3}s_1 + \tfrac{5}{3}s_2 = 8 - x_3 - 3x_4 - s_1 + s_2.
\]

Actualizar objetivo ($c_1=2$):
\[
  Z = \tfrac{32}{3} - \tfrac{1}{3}x_1 - \tfrac{1}{3}x_4 - \tfrac{1}{3}s_1 - \tfrac{1}{3}s_2\;+\;\tfrac{2}{3}\cdot\tfrac{2}{3}\cdot x_1\;\cdots
\]

Usando directamente $Z = 2x_1 + x_2 + 3x_3 + 2x_4$ con la nueva solución:
$x_1^* = 2$, $x_2^* = 8$, $x_3^*=x_4^*=0$. $Z^* = 2(2)+8 = \mathbf{12}$.

Nuevo diccionario:
\begin{align*}
x_1 &= 2 - \tfrac{3}{2}x_3 + \tfrac{1}{2}x_4 + \tfrac{1}{2}s_1 - s_2 \\
x_2 &= 8 - x_3 - 3x_4 - s_1 + s_2 \\
Z   &= 12 - x_3 - s_2
\end{align*}

Costos reducidos: $x_3: -1 \leq 0$; $x_4: 0$; $s_1: 0$; $s_2: -1 \leq 0$. Óptimo \textbf{[BT §3.2]}.

\item \textbf{$a_{13}$: $5 \to 3$.}

$x_3$ es básica. La columna cambia a $(3, 4)'$. Nueva base $B' = \bigl[\!\begin{smallmatrix}2&3\\1&4\end{smallmatrix}\!\bigr]$, $\det(B')=5$.

\[
  B'^{-1} = \frac{1}{5}\begin{pmatrix}4 & -3 \\ -1 & 2\end{pmatrix}.
\]

Nueva SBF \textbf{[BT §5.3]}:
\[
  x_{\mathcal{B}'} = B'^{-1}b = \frac{1}{5}\begin{pmatrix}4(20)-3(12)\\-1(20)+2(12)\end{pmatrix} = \frac{1}{5}\begin{pmatrix}44\\4\end{pmatrix} = \begin{pmatrix}8.8\\0.8\end{pmatrix} \geq 0\;\checkmark.
\]

Nuevos precios duales: $\mu' = (1,3)\cdot\frac{1}{5}\begin{pmatrix}4&-3\\-1&2\end{pmatrix} = \frac{1}{5}(1,3) = \left(\tfrac{1}{5},\tfrac{3}{5}\right)$.

Costos reducidos:
\begin{align*}
\bar{c}_{x_1}' &= 1 - \tfrac{1}{5}(2) - \tfrac{3}{5}(2) = 1 - \tfrac{2}{5} - \tfrac{6}{5} = -\tfrac{3}{5} \leq 0 \\
\bar{c}_{x_4}' &= 2 - \tfrac{1}{5}(5) - \tfrac{3}{5}(2) = 2-1-\tfrac{6}{5} = -\tfrac{1}{5} \leq 0 \\
\bar{c}_{s_1}' &= -\tfrac{1}{5} \leq 0;\quad \bar{c}_{s_2}' = -\tfrac{3}{5} \leq 0
\end{align*}

Base óptima. \textbf{Nueva solución:} $x_2^* = 8.8$, $x_3^* = 0.8$, $Z^* = 1(8.8) + 3(0.8) = 8.8 + 2.4 = \mathbf{11.2}$.

\end{enumerate}

% ============================================================
\section*{Problema 6 --- Dual de transporte con cotas superiores}
\addcontentsline{toc}{section}{Problema 6}

\textbf{Enunciado.} Construir el dual del problema de transporte con cotas superiores:
\[
\begin{array}{ll}
\min & \sum_{i,j} c_{ij} x_{ij} \\
\text{s.a.} & \sum_i x_{ij} \geq d_j, \quad \forall j \quad \text{(demandas)} \\
& \sum_j x_{ij} \leq o_i, \quad \forall i \quad \text{(ofertas)} \\
& x_{ij} \leq u_{ij}, \quad \forall i,j \quad \text{(cotas)} \\
& x_{ij} \geq 0
\end{array}
\]

\begin{marco}
\textbf{[BT §4.2]} Reglas de dualización para $\min$: restricción $\leq$ $\Rightarrow$ dual $\geq 0$; restricción $\geq$ $\Rightarrow$ dual $\leq 0$ (o reescribir como $-\sum_i x_{ij} \leq -d_j$ y usar dual $\geq 0$). El objetivo dual es $\max$ del producto de duales por RHS.
\end{marco}

\subsection*{Solución}

Reescribimos las restricciones $\geq$ como $\leq$: $-\sum_i x_{ij} \leq -d_j$ (duales $v_j \geq 0$).

Variables duales: $v_j \geq 0$ para demandas, $u_i \geq 0$ para ofertas, $w_{ij} \geq 0$ para cotas.

La restricción dual de $x_{ij}$ (variable $\geq 0$ en $\min$, tipo $\geq$):
\[
  (-1)\cdot v_j + (1)\cdot u_i + (1)\cdot w_{ij} \geq c_{ij} \;\Leftrightarrow\; u_i - v_j + w_{ij} \geq c_{ij}.
\]

Objetivo dual: $\max\; \bigl(-\sum_j d_j v_j + \sum_i o_i(-u_i) - \sum_{ij} u_{ij} w_{ij}\bigr)$.

Nota: para las restricciones de oferta, el signo del RHS en la forma $\leq$ es $+o_i$, y para las de cota es $+u_{ij}$. Por lo tanto:

\[
\boxed{
\begin{array}{ll}
\max & -\displaystyle\sum_j d_j v_j - \sum_i o_i u_i - \sum_{i,j} u_{ij} w_{ij} \\[6pt]
\text{s.a.} & u_i - v_j + w_{ij} \geq c_{ij}, \quad \forall i,j \\[2pt]
& v_j,\; u_i,\; w_{ij} \geq 0
\end{array}
}
\]

\textbf{Interpretación:} $v_j$ es el precio (negativo) de satisfacer la demanda $j$; $u_i$ es el costo de la oferta $i$; $w_{ij}$ es el costo de la cota superior de $x_{ij}$. La dualidad débil garantiza que $\text{dual} \leq \text{primal}$ para toda solución factible del par.

% ============================================================
\section*{Problema 7 --- Verdadero/Falso (4 afirmaciones)}
\addcontentsline{toc}{section}{Problema 7}

\textbf{Enunciado.} Comente sobre la veracidad de las siguientes afirmaciones.

\begin{marco}
\textbf{[BT §4.2 Teo.\ 4.2]} Dualidad fuerte: si el primal tiene solución óptima finita, el dual también; consecuencia: primal ilimitado $\Rightarrow$ dual infactible. \\[2pt]
\textbf{[BT §4.3 Teo.\ 4.5]} Holgura complementaria: $y_i^* > 0 \Rightarrow$ restricción primal activa; $x_j^* > 0 \Rightarrow$ restricción dual activa. \\[2pt]
\textbf{[BT §5.1]} Disminuir $c_j$ de una variable no básica mantiene $\bar{c}_j$ aún $\leq 0$. \\[2pt]
\textbf{[Ch §7]} En la Fase I del método de dos fases, la función auxiliar se trata como una maximización equivalente: $\min\sum a_k \equiv \max(-\sum a_k)$.
\end{marco}

\subsection*{Solución}

\begin{enumerate}[a)]

\item \textbf{``Si al iterar el simplex se obtiene una solución infactible, el dual es no acotado.''}

\textbf{FALSO.} El método simplex primal nunca produce soluciones infactibles: en cada iteración mantiene $x_{\mathcal{B}} = B^{-1}b \geq 0$ \textbf{[BT §3.1]}. La afirmación confunde la dirección correcta de la dualidad fuerte: primal ilimitado $\Rightarrow$ dual infactible \textbf{[BT §4.2 Teo.\ 4.2]}. La recíproca (dual infactible $\Rightarrow$ primal ilimitado) \emph{no} es siempre cierta: el primal puede ser infactible también.

\item \textbf{``Disminuir $c_j$ de una variable no básica nunca afecta la solución óptima.''}

\textbf{VERDADERO.} Sea $x_j$ no básica con $\bar{c}_j \leq 0$ (criterio de optimalidad). Nuevo costo reducido \textbf{[BT §5.1]}:
\[
  \bar{c}_j' = c_j' - \mu a_j = \bar{c}_j + \underbrace{(c_j' - c_j)}_{= -\delta,\;\delta\geq 0} \leq \bar{c}_j \leq 0.
\]
La variable $x_j$ sigue siendo no atractiva; la base, la SBF y $Z^*$ no cambian.

\item \textbf{``Si una variable dual óptima es positiva, la restricción primal correspondiente puede no ser de igualdad.''}

\textbf{FALSO.} Por holgura complementaria \textbf{[BT §4.3 Teo.\ 4.5]}:
\[
  y_i^*(b_i - a_i' x^*) = 0. \quad \text{Si } y_i^* > 0 \;\Rightarrow\; b_i - a_i'x^* = 0 \;\text{necesariamente.}
\]
Luego la restricción $i$ \emph{debe} ser activa en el óptimo primal.

\item \textbf{``En el método de dos fases para minimización, la función auxiliar de Fase I debe ser de maximización.''}

\textbf{VERDADERO.} La Fase I minimiza $\sum_k a_k$ (suma de variables artificiales), lo que es equivalente a maximizar $-\sum_k a_k$ \textbf{[Ch §7]}. Al aplicar el simplex de maximización se obtiene la misma base que al minimizar la función auxiliar. Si el óptimo de Fase I es $0$ (todas las artificiales fuera de la base), se tiene una SBF factible para el problema original.

\end{enumerate}

% ============================================================
\section*{Problema 8 --- Restricción inactiva y efecto en $B^{-1}b$}
\addcontentsline{toc}{section}{Problema 8}

\textbf{Enunciado.} Si se aumenta el RHS de una restricción inactiva en la solución óptima, ¿cómo cambia $B^{-1}b$? Justifique.

\begin{marco}
\textbf{[BT §3.1]} Los valores de la SBF son $x_{\mathcal{B}} = B^{-1}b$. Si $b \to b + \Delta e_i$, entonces $x_{\mathcal{B}}' = x_{\mathcal{B}} + \Delta B^{-1}e_i$. \\[2pt]
\textbf{[BT §4.3 Teo.\ 4.5]} Restricción inactiva $\Rightarrow$ su holgura $s_i > 0 \Rightarrow s_i$ básica $\Rightarrow y_i^* = 0$. \\[2pt]
\textbf{[BT §5.3]} $B^{-1}e_i = e_k$ donde $k$ es la posición de $s_i$ en la base $\mathcal{B}$.
\end{marco}

\subsection*{Solución}

Sea la restricción $i$ inactiva en el óptimo: $s_i = b_i - a_i'x^* > 0$, por lo que $s_i$ es variable básica en alguna posición $k$ de $\mathcal{B}$. La columna de $s_i$ en la matriz ampliada es $e_i$ (vector canónico $i$). Su representación en la base actual es $B^{-1}e_i = e_k$ (el vector canónico $k$-ésimo de $\mathbb{R}^m$) \textbf{[BT §5.3]}.

Al aumentar $b_i$ en $\Delta > 0$:
\[
  x_{\mathcal{B}}' = x_{\mathcal{B}} + \Delta \cdot B^{-1}e_i = x_{\mathcal{B}} + \Delta e_k.
\]

\textbf{Consecuencias:}
\begin{enumerate}[(i)]
  \item Solo la componente $k$-ésima cambia: $s_i' = s_i + \Delta > 0$. Todas las demás variables básicas (incluidas las variables de decisión) permanecen iguales.
  \item La base sigue siendo factible: $s_i' = s_i + \Delta > 0$ \checkmark.
  \item Los costos reducidos no cambian (dependen de $\mu = c_{\mathcal{B}}B^{-1}$, no de $b$).
  \item $Z^*$ no cambia: $\Delta Z = y_i^* \cdot \Delta = 0$, pues restricción inactiva $\Rightarrow y_i^* = 0$ por holgura complementaria \textbf{[BT §4.3 Teo.\ 4.5]}.
\end{enumerate}

En resumen: aumentar el RHS de una restricción inactiva solo incrementa la holgura de esa restricción en la misma cantidad, sin alterar ninguna variable de decisión ni el valor óptimo.

% ============================================================
\section*{Problema 9 --- Eliminar restricción 3 de P3}
\addcontentsline{toc}{section}{Problema 9}

\textbf{Enunciado.} Retome el problema agrícola (P3) sin la restricción $x_3 \leq 350$. Encuentre la nueva solución óptima sin resolver desde cero.

\begin{marco}
\textbf{[BT §4.3 Teo.\ 4.5]} $y_3^* = 2.5 > 0 \Rightarrow$ la restricción 3 era activa en el óptimo original ($x_3^* = 350$). Eliminarla cambia la región factible, mejorando (o dejando igual) el óptimo. \\[2pt]
\textbf{[BT §5.3]} Al eliminar una restricción activa cuyo dual es positivo, la base actual ya no es óptima para el nuevo problema y el objetivo mejora.
\end{marco}

\subsection*{Solución}

$y_3^* = 2.5 > 0$ implica que la restricción $x_3 \leq 350$ era activa en el óptimo original \textbf{[BT §4.3 Teo.\ 4.5]}. Al eliminarla, $x_3$ puede tomar valores mayores que 350.

Sin restricción 3, el problema es:
\[
\begin{array}{ll}
\max & 150x_1 + 200x_2 - 10x_3 \\
\text{s.a.} & x_1 + x_2 \leq 45 \\
& 6x_1 + 10x_2 - x_3 \leq 0 \\
& 5x_1 \leq 140 \;\Rightarrow\; x_1 \leq 28 \\
& 4x_2 \leq 120 \;\Rightarrow\; x_2 \leq 30 \\
& x_1, x_2, x_3 \geq 0
\end{array}
\]

Dado que $-10 < 0$, el objetivo se maximiza minimizando $x_3$. La restricción (2) actúa como cota inferior: $x_3 \geq 6x_1 + 10x_2$. El mínimo factible es $x_3^* = 6x_1 + 10x_2$.

Sustituyendo en el objetivo:
\[
  Z = 150x_1 + 200x_2 - 10(6x_1 + 10x_2) = 90x_1 + 100x_2.
\]

Problema reducido:
\[
\max\; 90x_1 + 100x_2 \quad \text{s.a.}\quad x_1+x_2\leq 45,\; x_1\leq 28,\; x_2\leq 30,\; x_1,x_2\geq 0.
\]

Evaluamos los vértices de la región factible:
\begin{center}
\begin{tabular}{cccc}
\hline
Vértice & $x_1$ & $x_2$ & $Z$\\
\hline
$O$ & 0 & 0 & 0 \\
$A$ & 15 & 30 & $90(15)+100(30)=1350+3000=4350$ \\
$B$ & 28 & 17 & $90(28)+100(17)=2520+1700=4220$ \\
$C$ & 28 & 0 & $2520$ \\
$D$ & 0 & 30 & $3000$ \\
\hline
\end{tabular}
\end{center}

\textbf{Óptimo en $A$:} $x_1^* = 15$, $x_2^* = 30$, $x_3^* = 6(15)+10(30) = 90+300 = 390$,
$Z^* = 150(15)+200(30)-10(390) = 2250+6000-3900 = \mathbf{4350}$ M\$.

% ============================================================
\section*{Problema 10 --- Eliminar variables de P5}
\addcontentsline{toc}{section}{Problema 10}

\textbf{Enunciado.} Para el problema P5: (a) resolver sin la variable $x_4$; (b) resolver sin la variable $x_3$.

\begin{marco}
\textbf{[BT §5.2]} Variable no básica con $\bar{c}_j < 0$: eliminarla no cambia la base ni la solución óptima. \\[2pt]
\textbf{[BT §5.2]} Variable básica: eliminarla modifica el modelo; se debe re-optimizar partiendo del diccionario óptimo del nuevo problema.
\end{marco}

\subsection*{Solución}

\begin{enumerate}[a)]

\item \textbf{Sin $x_4$.}

En el diccionario óptimo de P5, $x_4$ es no básica con $\bar{c}_{x_4} = -\tfrac{1}{3} < 0$ \textbf{[BT §5.2]}. Eliminar $x_4$ no afecta la base $\{x_2, x_3\}$ ni la SBF.

\textbf{Solución sin cambios:} $x_2^* = 20/3$, $x_3^* = 4/3$, $Z^* = 32/3$.

\item \textbf{Sin $x_3$.}

Nuevo problema:
\[
\begin{array}{ll}
\max & x_1 + x_2 + 2x_4 \\
\text{s.a.} & 2x_1 + 2x_2 + 5x_4 \leq 20 \\
& 2x_1 + x_2 + 2x_4 \leq 12 \\
& x_1, x_2, x_4 \geq 0
\end{array}
\]

\textit{Diccionario inicial:} $s_1 = 20 - 2x_1 - 2x_2 - 5x_4$; $s_2 = 12-2x_1-x_2-2x_4$; $Z = x_1+x_2+2x_4$.

\medskip
\textbf{Iteración 1.} Entra $x_4$ ($\bar{c}_{x_4} = 2$). Razones: $20/5=4$ (fila $s_1$), $12/2=6$ (fila $s_2$). Sale $s_1$.

Pivote: $x_4 = 4 - \tfrac{2}{5}x_1 - \tfrac{2}{5}x_2 - \tfrac{1}{5}s_1$.

Actualizar:
\begin{align*}
s_2 &= 12 - 2x_1 - x_2 - 2\bigl(4 - \tfrac{2}{5}x_1 - \tfrac{2}{5}x_2 - \tfrac{1}{5}s_1\bigr) = 4 - \tfrac{6}{5}x_1 - \tfrac{1}{5}x_2 + \tfrac{2}{5}s_1 \\
Z   &= x_1 + x_2 + 2\bigl(4 - \tfrac{2}{5}x_1 - \tfrac{2}{5}x_2 - \tfrac{1}{5}s_1\bigr) = 8 + \tfrac{1}{5}x_1 + \tfrac{1}{5}x_2 - \tfrac{2}{5}s_1
\end{align*}

\medskip
\textbf{Iteración 2.} Entra $x_1$ ($\bar{c}_{x_1} = \tfrac{1}{5} > 0$). Razones: fila $x_4$: $4/(2/5)=10$; fila $s_2$: $4/(6/5)=10/3$. Sale $s_2$ ($\theta^* = 10/3$).

Pivote: $x_1 = \tfrac{10}{3} - \tfrac{1}{6}x_2 + \tfrac{1}{3}s_1 - \tfrac{5}{6}s_2$.

Actualizar:
\begin{align*}
x_4 &= 4 - \tfrac{2}{5}\bigl(\tfrac{10}{3}-\tfrac{1}{6}x_2+\tfrac{1}{3}s_1-\tfrac{5}{6}s_2\bigr) - \tfrac{2}{5}x_2 - \tfrac{1}{5}s_1 \\
    &= 4 - \tfrac{4}{3} + \tfrac{1}{15}x_2 - \tfrac{2}{15}s_1 + \tfrac{1}{3}s_2 - \tfrac{2}{5}x_2 - \tfrac{1}{5}s_1 \\
    &= \tfrac{8}{3} - \tfrac{1}{3}x_2 - \tfrac{1}{3}s_1 + \tfrac{1}{3}s_2 \\
Z   &= 8 + \tfrac{1}{5}\bigl(\tfrac{10}{3}-\tfrac{1}{6}x_2+\tfrac{1}{3}s_1-\tfrac{5}{6}s_2\bigr) + \tfrac{1}{5}x_2 - \tfrac{2}{5}s_1 \\
    &= 8 + \tfrac{2}{3} + \tfrac{1}{6}x_2 - \tfrac{1}{6}s_2 - \tfrac{1}{3}s_1 \\
    &= \tfrac{26}{3} + \tfrac{1}{6}x_2 - \tfrac{1}{3}s_1 - \tfrac{1}{6}s_2
\end{align*}

\medskip
\textbf{Iteración 3.} Entra $x_2$ ($\bar{c}_{x_2} = \tfrac{1}{6} > 0$). Razones: fila $x_1$: $(10/3)/(1/6)=20$; fila $x_4$: $(8/3)/(1/3)=8$. Sale $x_4$ ($\theta^* = 8$).

Pivote: $x_2 = 8 - 3x_4 - s_1 + s_2$.

Actualizar:
\begin{align*}
x_1 &= \tfrac{10}{3} - \tfrac{1}{6}(8-3x_4-s_1+s_2) + \tfrac{1}{3}s_1 - \tfrac{5}{6}s_2 = 2 + \tfrac{1}{2}x_4 + \tfrac{1}{2}s_1 - s_2 \\
Z   &= \tfrac{26}{3} + \tfrac{1}{6}(8-3x_4-s_1+s_2) - \tfrac{1}{3}s_1 - \tfrac{1}{6}s_2 = 10 - \tfrac{1}{2}x_4 - \tfrac{1}{2}s_1
\end{align*}

Costos reducidos: $x_4: -\tfrac{1}{2}\leq 0$; $s_1: -\tfrac{1}{2}\leq 0$; $s_2: 0$. Óptimo \textbf{[BT §3.2]}.

\textbf{Nueva solución:} $x_1^* = 2$, $x_2^* = 8$, $x_4^* = 0$, $Z^* = \mathbf{10}$.

\end{enumerate}

% ============================================================
\section*{Problema 11 --- Degeneración primal $\leftrightarrow$ dual}
\addcontentsline{toc}{section}{Problema 11}

\textbf{Enunciado.} ¿Cómo se refleja la degeneración primal óptima en el dual? ¿Es cierta la situación inversa?

\begin{marco}
\textbf{[BT §4.3 Teo.\ 4.5]} Holgura complementaria: $y_i^*(b_i - a_i'x^*) = 0$ y $x_j^*((A'y^*)_j - c_j) = 0$. \\[2pt]
Degeneración primal: $x_{Bk}^* = 0$ para algún $k \in \mathcal{B}$. Degeneración dual: $y_{Bi}^* = 0$ para alguna variable básica del dual.
\end{marco}

\subsection*{Solución}

\textbf{Degeneración primal $\Rightarrow$ múltiples soluciones óptimas duales.}

Sea $x_{Bk}^* = 0$ para algún $k \in \mathcal{B}$. El mismo punto extremo primal puede ser representado por más de una base (varias bases degeneradas en el mismo vértice). Cada base óptima $\mathcal{B}^{(r)}$ genera precios duales distintos:
\[
  y^{(r)} = \bigl(B^{(r)}\bigr)^{-T} c_{\mathcal{B}^{(r)}}.
\]
Todos son óptimos duales con $W = Z^*$ (por dualidad fuerte \textbf{[BT §4.2 Teo.\ 4.2]}). Por lo tanto, el dual tiene \textbf{múltiples soluciones óptimas}.

\textbf{Situación inversa: degeneración dual $\Rightarrow$ múltiples soluciones óptimas primales.}

Sea $y_{Bi}^* = 0$ para alguna variable básica del dual. La condición de holgura complementaria $y_i^*(b_i - a_i'x^*) = 0$ se satisface trivialmente sin imponer que la restricción primal $i$ sea activa. Esto significa que $x^*$ puede satisfacer con holgura la restricción $i$ (distinta de cero), y aun así ser óptimo. Otro punto $\hat{x}^*$ con distinta holgura en la restricción $i$ también puede ser óptimo. Por tanto el primal tiene \textbf{múltiples soluciones óptimas}.

\medskip
\noindent\textbf{Resumen:}
\begin{center}
\begin{tabular}{lcl}
\hline
Degeneración primal ($x_{Bk}^* = 0$) & $\Leftrightarrow$ & Dual tiene múltiples soluciones óptimas \\
Degeneración dual ($y_{Bi}^* = 0$) & $\Leftrightarrow$ & Primal tiene múltiples soluciones óptimas \\
\hline
\end{tabular}
\end{center}

% ============================================================
\section*{Problema 12 --- Dual infactible para qué valores de $c$}
\addcontentsline{toc}{section}{Problema 12}

\textbf{Enunciado.} Dado el problema $\max\; cx_1 - 4x_2$ s.a.\ $x_1 - x_2 \leq 1$, $x_1, x_2 \geq 0$: ¿para qué valores de $c$ el dual es infactible?

\begin{marco}
\textbf{[BT §4.2 Teo.\ 4.2]} Primal ilimitado $\Rightarrow$ dual infactible. Dual factible $\Rightarrow$ primal acotado (dualidad débil). \\[2pt]
\textbf{[BT §4.1 Teo.\ 4.1]} Dualidad débil: para toda solución primal factible $\hat{x}$ y dual factible $\hat{y}$, $c\hat{x} \leq \hat{y}b$.
\end{marco}

\subsection*{Solución}

\textbf{Construcción del dual.}

Primal: $\max\; cx_1 - 4x_2$; $x_1 - x_2 \leq 1$; $x_1, x_2 \geq 0$.

Dual: $\min\; y$; s.a.\ $y \geq c$ (de $x_1$); $-y \geq -4$, i.e., $y \leq 4$ (de $x_2$); $y \geq 0$.

El dual es factible ssi existe $y$ con $c \leq y \leq 4$ y $y \geq 0$, lo que requiere $c \leq 4$ (y $c \leq 4$, $4 \geq 0$, ambas compatibles cuando $c \leq 4$).

\textbf{El dual es infactible $\Leftrightarrow c > 4$.}

\textit{Verificación directa \textbf{[BT §4.2 Teo.\ 4.2]}:} si $c > 4$, tome la solución $x_1 = t,\; x_2 = t-1$ para $t \geq 1$: la restricción $x_1 - x_2 = 1 \leq 1$ es factible, y $Z = (c-4)t + 4 \to +\infty$ cuando $t \to \infty$. El primal es ilimitado, luego el dual es infactible. $\square$

% ============================================================
\section*{Problema 13 --- Probar $cx^* \leq y^*d$}
\addcontentsline{toc}{section}{Problema 13}

\textbf{Enunciado.} Sea $y^*$ óptimo del dual del problema primal con RHS $b$. Sea $x^*$ óptimo del primal con RHS $d$ (no necesariamente igual a $b$). Probar que $cx^* \leq y^*d$.

\begin{marco}
\textbf{[BT §4.1 Teo.\ 4.1]} Dualidad débil: para todo $\hat{x}$ factible en el primal y $\hat{y}$ factible en el dual, $c\hat{x} \leq \hat{y}b$. Equivalentemente: si $A\hat{x} \leq b$, $\hat{x} \geq 0$ y $A'\hat{y} \geq c$, $\hat{y} \geq 0$, entonces $c\hat{x} \leq (A\hat{x})'\hat{y} \leq b'\hat{y}$.
\end{marco}

\subsection*{Solución}

\textbf{Hipótesis:}
\begin{itemize}
  \item $y^*$ factible para el dual original: $A'y^* \geq c$, $y^* \geq 0$.
  \item $x^*$ factible para el primal modificado: $Ax^* \leq d$, $x^* \geq 0$.
\end{itemize}

\textbf{Paso 1.} Como $A'y^* \geq c$ y $x^* \geq 0$:
\[
  (A'y^*)'x^* \geq c'x^* \;\Leftrightarrow\; (y^*)'(Ax^*) \geq cx^*.
\]

\textbf{Paso 2.} Como $Ax^* \leq d$ y $y^* \geq 0$:
\[
  (y^*)'(Ax^*) \leq (y^*)'d = y^*d.
\]

\textbf{Combinando:}
\[
  \boxed{cx^* \leq (y^*)'(Ax^*) \leq y^*d.} \qquad \square
\]

\textbf{Interpretación:} los precios sombra $y^*$ del sistema original acotan el valor óptimo del primal incluso cuando el RHS se modifica a $d \neq b$.

% ============================================================
\section*{Problema 14 --- Verdadero/Falso (5 afirmaciones)}
\addcontentsline{toc}{section}{Problema 14}

\textbf{Enunciado.} Demuestre o desmiente las siguientes afirmaciones.

\begin{marco}
\textbf{[BT §3.3]} $\Delta Z = \bar{c}_s \cdot \theta^*$. Mayor $\bar{c}_s$ no garantiza mayor $\Delta Z$ pues $\theta^*$ puede ser pequeño. \\[2pt]
\textbf{[BT §3.3]} Degeneración ($\theta^* = 0$): posible pero no obligatorio. \\[2pt]
\textbf{[BT §2.1]} El conjunto de soluciones óptimas de un PL es convexo (y poliedrico), y si contiene dos puntos distintos contiene el segmento entre ellos: infinitos puntos. \\[2pt]
\textbf{[BT §4.3 Teo.\ 4.5]} $y_i^* > 0 \Rightarrow$ restricción activa. La recíproca es falsa (degeneración dual). \\[2pt]
\textbf{[BT §3.2]} Criterio de optimalidad: todos los costos reducidos $\leq 0$.
\end{marco}

\subsection*{Solución}

\begin{enumerate}[a)]

\item \textbf{``Elegir la variable con mayor costo reducido maximiza el aumento de $Z$ en esa iteración.''}

\textbf{FALSO.} El cambio en el objetivo es $\Delta Z = \bar{c}_s \cdot \theta^*$, donde $\theta^* = \min_i \{b_i / \bar{a}_{is} : \bar{a}_{is} > 0\}$ depende de la columna $s$ seleccionada \textbf{[BT §3.3]}. Una variable con mayor $\bar{c}_s$ puede tener $\theta^*$ muy pequeño.

\textit{Contraejemplo:} $\bar{c}_1 = 10,\; \theta_1^* = 0.1 \Rightarrow \Delta Z_1 = 1$; $\bar{c}_2 = 3,\; \theta_2^* = 2 \Rightarrow \Delta Z_2 = 6$. Elegir $x_2$ da mayor incremento.

\item \textbf{``Si ambos extremos del diccionario son degenerados, $Z$ puede no aumentar.''}

\textbf{CORRECTO (con matiz).} La degeneración solo garantiza $\Delta Z = 0$ cuando $\theta^* = 0$ (paso degenerado \textbf{[BT §3.3]}). Si la variable básica con valor 0 no es la que determina la razón mínima, puede ocurrir $\theta^* > 0$ y $\Delta Z > 0$. Sin embargo, \emph{si} ambas bases son degeneradas en el sentido de que la razón mínima resulta 0 en la iteración, efectivamente $\Delta Z = 0$.

\item \textbf{``Si un PL tiene múltiples soluciones óptimas, tiene un número incontable de ellas.''}

\textbf{VERDADERO.} Sean $x^*$ y $\hat{x}$ dos soluciones óptimas distintas con $cx^* = c\hat{x} = Z^*$. Para todo $\lambda \in [0,1]$:
\[
  \tilde{x} = \lambda x^* + (1-\lambda)\hat{x}.
\]
Factibilidad: $A\tilde{x} \leq b$ (por convexidad del politopo \textbf{[BT §2.1]}), $\tilde{x} \geq 0$. Optimalidad: $c\tilde{x} = \lambda Z^* + (1-\lambda)Z^* = Z^*$. Como $\lambda$ varía continuamente en $[0,1]$, hay incontables soluciones óptimas. $\square$

\item \textbf{``Si una restricción es activa en el óptimo primal, la variable dual correspondiente es positiva.''}

\textbf{FALSO.} La holgura complementaria \textbf{[BT §4.3 Teo.\ 4.5]} establece $y_i^*(b_i - a_i'x^*) = 0$. Si $b_i - a_i'x^* = 0$ (restricción activa), la ecuación se satisface para \emph{cualquier} $y_i^* \geq 0$, incluyendo $y_i^* = 0$ (degeneración dual). La implicación válida es $y_i^* > 0 \Rightarrow$ restricción activa, no su recíproco.

\item \textbf{``En un diccionario óptimo factible ninguna variable no básica tiene costo reducido positivo.''}

\textbf{VERDADERO.} Por contradicción: si $\bar{c}_j > 0$ para alguna no básica $x_j$, aumentar $x_j$ incrementa $Z$ \textbf{[BT §3.2]}. Si $\theta^* > 0$, existe una SBF factible con $Z' = Z + \bar{c}_j\theta^* > Z^*$, contradiciendo la optimalidad. Si $\theta^* = 0$, el paso degenerado eventualmente podría llevar a $Z' > Z^*$ o revelar no-acotado. En ambos casos el diccionario no puede ser óptimo. $\square$

\end{enumerate}

% ============================================================
\section*{Problema 15 --- Chvátal Teorema 3.2 (Regla lexicográfica)}
\addcontentsline{toc}{section}{Problema 15}

\textbf{Enunciado.} Lea la prueba del Teorema 3.2 de Chvátal. Explique cómo la expresión (3.9) (razón lexicográfica) se usa \emph{solo} cuando hay empate en la variable saliente, y por qué esto es suficiente para garantizar la no ciclación.

\begin{marco}
\textbf{[Ch §3 Teo.\ 3.2]} Regla lexicográfica: si en cada iteración se elige la variable saliente por la razón lexicográfica mínima, el método simplex termina en un número finito de iteraciones. \\[2pt]
La idea de la prueba: el vector fila del objetivo (aumentado con las columnas de la identidad) crece estrictamente en sentido lexicográfico en cada iteración, y hay un número finito de bases distintas, por lo que no puede haber ciclos.
\end{marco}

\subsection*{Solución}

\textbf{Expresión (3.9):} para la variable entrante $x_s$ y cada fila $i$ con $\bar{a}_{is} > 0$, la razón lexicográfica es:
\[
  r_i = \left(\frac{\bar{b}_i}{\bar{a}_{is}},\; \frac{\bar{a}_{i1}}{\bar{a}_{is}},\; \frac{\bar{a}_{i2}}{\bar{a}_{is}},\; \ldots,\; \frac{\bar{a}_{im}}{\bar{a}_{is}}\right) \in \mathbb{R}^{m+1}.
\]
La variable saliente es $x_{B(t)}$ donde $t = \arg\min_{\text{lex}} r_i$ entre los $i$ con $\bar{a}_{is} > 0$.

\textbf{Implementación selectiva:}

\begin{enumerate}[(1)]
  \item Calcular $\theta^* = \min_{i:\,\bar{a}_{is}>0} \bar{b}_i / \bar{a}_{is}$ (razón mínima ordinaria).
  \item Si una sola fila $i$ alcanza $\theta^*$: usar esa fila como variable saliente. \emph{La expresión (3.9) no se necesita.}
  \item Si varias filas $T = \{i : \bar{b}_i/\bar{a}_{is} = \theta^*\}$ empatan: aplicar la expresión (3.9) solo sobre $T$, comparando lexicográficamente las columnas $1, 2, \ldots, m$ hasta desempatar.
\end{enumerate}

\textbf{¿Por qué es suficiente?}

La prueba del Teorema 3.2 en Chvátal demuestra que el vector fila del objetivo aumentado $(\bar{Z}, \bar{a}_{01}, \ldots, \bar{a}_{0m})$ es \emph{lexicográficamente positivo} en todo momento, y que crece estrictamente en cada iteración (incluso cuando $\theta^* = 0$, es decir, en pasos degenerados). Este crecimiento lexicográfico estricto implica que ninguna base puede repetirse, garantizando terminación finita.

Cuando no hay empate en $\theta^*$, el paso no es degenerado ($\theta^* > 0$) y el objetivo $Z$ aumenta estrictamente. Esto por sí solo impide la repetición de bases. Solo en pasos \emph{potencialmente} degenerados (varios $i$ con $\bar{b}_i = 0$ empatan en $\theta^* = 0$) es necesario recurrir a (3.9) para garantizar el crecimiento lexicográfico del objetivo aumentado y así evitar ciclos. Por ello, aplicar la regla lexicográfica solo al desempatar es suficiente y, en la práctica, ocurre raramente.

% ============================================================
\section*{Problema 16 --- Chvátal 5.9 y 5.10}
\addcontentsline{toc}{section}{Problema 16}

\textbf{Enunciado.} (a) Chvátal ejercicio 5.9: probar que si $x^*$ es primal factible, $y^*$ es dual factible y $cx^* = y^*b$, entonces ambos son óptimos. (b) Chvátal ejercicio 5.10: probar que si $x^*$ es la única SBF óptima y no degenerada, el dual tiene solución óptima única.

\begin{marco}
\textbf{[Ch §5 Teo.\ 5.1 / BT §4.1 Teo.\ 4.1]} Dualidad débil: $cx \leq yb$ para todo $x$ primal factible e $y$ dual factible. \\[2pt]
\textbf{[Ch §5 Teo.\ 5.3 / BT §4.2 Teo.\ 4.2]} Dualidad fuerte: si el primal tiene óptimo finito, el dual también y los valores óptimos coinciden. \\[2pt]
\textbf{[BT §4.3 Teo.\ 4.5]} Holgura complementaria: $(x^*, y^*)$ par óptimo $\Leftrightarrow$ se satisfacen las condiciones de complementariedad.
\end{marco}

\subsection*{Solución}

\begin{enumerate}[a)]

\item \textbf{Chvátal 5.9: si $cx^* = y^*b$, entonces $x^*$ y $y^*$ son óptimos.}

\textbf{Demostración.}

\textit{$x^*$ es primal óptimo:} Para cualquier $x$ primal factible, por dualidad débil \textbf{[Ch §5 Teo.\ 5.1]}:
\[
  cx \leq y^*b = cx^*.
\]
Luego $cx \leq cx^*$ para todo $x$ factible, es decir, $x^*$ maximiza el primal. $\square$

\textit{$y^*$ es dual óptimo:} Para cualquier $y$ dual factible, por dualidad débil:
\[
  cx^* \leq yb \;\Rightarrow\; y^*b = cx^* \leq yb.
\]
Luego $y^*b \leq yb$ para todo $y$ dual factible, es decir, $y^*$ minimiza el dual. $\square$

La igualdad $cx^* = y^*b$ equivale a que se satisfagan las condiciones de holgura complementaria \textbf{[BT §4.3 Teo.\ 4.5]}:
\[
  y_i^*(b_i - a_i'x^*) = 0 \quad \forall i, \qquad x_j^*((A'y^*)_j - c_j) = 0 \quad \forall j.
\]

\item \textbf{Chvátal 5.10: si $x^*$ es la única SBF óptima no degenerada, el dual tiene solución óptima única.}

\textbf{Demostración.}

Sea $\mathcal{B}$ la base óptima asociada a $x^*$. Como $x^*$ es no degenerada: $x_{Bk}^* > 0$ para todo $k \in \mathcal{B}$.

Por holgura complementaria \textbf{[BT §4.3 Teo.\ 4.5]}: si $x_j^* > 0$ entonces la restricción dual $j$ es activa, es decir, $(A'y^*)_j = c_j$. Como \emph{todos} los $x_j^*$ para $j \in \mathcal{B}$ son estrictamente positivos, la $j$-ésima restricción dual es activa para todo $j \in \mathcal{B}$:
\[
  (A'y^*)_j = c_j \quad \forall j \in \mathcal{B} \;\Leftrightarrow\; B'y^* = c_{\mathcal{B}}.
\]

Como $B$ es no singular (es una base), el sistema $B'y^* = c_{\mathcal{B}}$ tiene solución única:
\[
  y^* = (B')^{-1}c_{\mathcal{B}} = B^{-T}c_{\mathcal{B}}.
\]

Toda solución dual óptima debe satisfacer esta ecuación (pues la no-degeneración primal fija las restricciones duales activas de todas las variables básicas). Por tanto \textbf{la solución dual óptima es única.} $\square$

\textbf{Interpretación geométrica:} cada variable básica estrictamente positiva ``activa'' exactamente su restricción dual, dejando el sistema $B'y^* = c_{\mathcal{B}}$ completamente determinado. La unicidad dual refleja que el vértice primal no degenerado está bien definido y no da lugar a ambigüedad en los precios sombra.

\end{enumerate}

\end{document}
